\begin{table}[htbp]
  \centering
  \caption{Minimalna długość odcinka $L_\mathrm{min}$ do osiągnięcia prędkości zadanej (scenariusz bazowy: $m=600$~t, $i=0$\textperthousand), przy maksymalnej mocy systemu. Pułap prędkości: AC $\approx$364~km/h, DC $\approx$332~km/h.}
  \label{tab:dlugosc-minimalna}
  \begin{tabular}{S[table-format=3.0] S[table-format=3.1] S[table-format=3.1] S[table-format=2.1]}
    \toprule
    {$v_\mathrm{zad}$ [km/h]} & {$L_\mathrm{min}$ AC [km]} & {$L_\mathrm{min}$ DC [km]} & {$\Delta$ (DC$-$AC) [km]} \\
    \midrule
    160 & 7.1 & 8.0 & 0.8 \\
    200 & 9.2 & 11.1 & 1.9 \\
    250 & 13.8 & 18.6 & 4.8 \\
    280 & 18.3 & 26.9 & 8.5 \\
    300 & 22.5 & 35.8 & 13.3 \\
    320 & 28.7 & 54.6 & 25.9 \\
    350 & 50.5 & {—} & {—} \\
    \bottomrule
  \end{tabular}
  \vspace{2pt}{\footnotesize\\ Symbol „---'' oznacza prędkość powyżej pułapu danego systemu (nieosiągalną na żadnej długości odcinka).}
\end{table}
